\chapterpage{Hashing and counting}{Eight problems. If a question asks
\emph{how many}, \emph{have I seen}, or \emph{group by}, the answer starts
with a \texttt{HashMap}. The prefix trick in P18 and P21 is the highest-value
pattern in this book.}

% ------------------------------------------------------------------
\problem{Count pairs summing to a target}{Easy}

Count the pairs of indices $i<j$ with \texttt{a[i] + a[j] == target}.

\idea{Walk once. For each element, the number of earlier elements that complete
the pair is already in the map. Count first, then insert.}

\begin{dia}
  \cellrowst{a}{0}{dead/1,dead/5,dead/7,hot/1}
  \idxrow{-0.62}{3}
  \node[ptr] at (2.4,0.78) {x = 1};
  \draw[flow] (2.4,0.62) -- (2.4,0.44);
  \node[note, anchor=west] at (3.3,0.35) {target 6, so look up $6-1=5$};
  \node[note, anchor=west] at (3.3,-0.2) {map says \texttt{5} seen once};
  \node[note, anchor=west] at (3.3,-0.75) {$\Rightarrow$ one new pair};
  \node[soft, minimum width=11mm] at (0.0,-1.8) {1:1};
  \node[soft, minimum width=11mm] at (1.2,-1.8) {5:1};
  \node[soft, minimum width=11mm] at (2.4,-1.8) {7:1};
  \node[note, anchor=east] at (-0.7,-1.8) {map before};
\end{dia}

\textbf{How it works.} Every pair has a later element and an earlier one. If
you always count from the point of view of the later element, each pair is
counted exactly once --- no division by two, no double counting.

That is why the order of the two statements matters. Looking up the
complement \emph{before} inserting the current value means the map holds only
elements strictly to the left. Insert first and an element can pair with
itself: a single \texttt{3} with target 6 would report a pair that does not
exist.

\solheading
\begin{lstlisting}[style=java]
public static long countPairs(int[] a, int target) {
    if (a == null) return 0;
    Map<Integer,Integer> seen = new HashMap<>();
    long pairs = 0;
    for (int x : a) {
        pairs += seen.getOrDefault(target - x, 0);
        seen.merge(x, 1, Integer::sum);
    }
    return pairs;
}
\end{lstlisting}

\cost{$O(n)$ time, $O(n)$ space.}
\trap{Returning \texttt{int}. With $200\,000$ identical elements the count is
about $2\times10^{10}$, which overflows.}

% ------------------------------------------------------------------
\problem{Longest subarray with equal counts}{Medium}

Return the length of the longest contiguous subarray holding as many even
numbers as odd ones.

\idea{Turn it into a running balance: $+1$ for odd, $-1$ for even. A subarray is
balanced exactly when the balance at its two ends is equal.}

\begin{dia}
  \draw[edge] (0,0) -- (5.4,0);
  \node[idx, anchor=east] at (-0.15,0) {0};
  \node[idx, anchor=east] at (-0.15,0.55) {1};
  \node[idx, anchor=east] at (-0.15,1.1) {2};
  \draw[edge, densely dotted] (0,0.55) -- (5.4,0.55);
  \draw[flow, draw=accent, line width=1pt]
    (0,0) -- (0.6,0.55) -- (1.2,0) -- (1.8,0.55) -- (2.4,0) -- (3.0,0.55);
  \foreach \x/\y in {0/0, 0.6/0.55, 1.2/0, 1.8/0.55, 2.4/0, 3.0/0.55} {
    \fill[accent] (\x,\y) circle (1.6pt);}
  \node[idx] at (0,-0.35) {$-1$};
  \node[idx] at (0.6,-0.35) {0};
  \node[idx] at (1.2,-0.35) {1};
  \node[idx] at (1.8,-0.35) {2};
  \node[idx] at (2.4,-0.35) {3};
  \node[idx] at (3.0,-0.35) {4};
  \node[note, anchor=west] at (3.4,0.55) {index};
  \draw[span] (0,-0.85) -- (2.4,-0.85);
  \caplabel{1.2}{-1.35}{same balance at both ends $\Rightarrow$ length 4};
  \node[note, anchor=west] at (3.4,0) {balance};
\end{dia}

\textbf{How it works.} Write $B_i$ for the balance after the first $i$
elements. The subarray from $j+1$ to $i$ has as many odds as evens exactly
when $B_i = B_j$, because the difference $B_i - B_j$ counts (odds $-$ evens)
inside that window.

So the problem becomes: for each balance value, how far apart are its two
furthest occurrences? Since the balance is scanned left to right, you only
ever need the \emph{earliest} index at which each value appeared --- any later
repeat measured against it gives the longest span. Hence \texttt{putIfAbsent}
semantics: write the index only when the value is new.

The seed \texttt{firstIndex.put(0, -1)} represents the empty prefix. Without
it, a subarray that is balanced from index 0 has no earlier index to measure
against and is missed entirely.

\solheading
\begin{lstlisting}[style=java]
public static int longestBalanced(int[] nums) {
    if (nums == null || nums.length == 0) return 0;
    Map<Integer,Integer> firstIndex = new HashMap<>();
    firstIndex.put(0, -1);                 // the empty prefix
    int balance = 0, best = 0;
    for (int i = 0; i < nums.length; i++) {
        balance += (nums[i] % 2 != 0) ? 1 : -1;
        Integer seen = firstIndex.get(balance);
        if (seen == null) firstIndex.put(balance, i);
        else best = Math.max(best, i - seen);
    }
    return best;
}
\end{lstlisting}

\cost{$O(n)$ time, $O(n)$ space.}
\trap{Three in one problem. Overwriting the stored index on a repeat measures
the shortest span. Omitting the \texttt{0 -> -1} seed loses answers starting
at index 0. And \texttt{nums[i] \% 2 == 1} is false for $-3$.}
\alsoknown{Task 4 of the assessment; and with 0 and 1 in place of even and
odd, the classic ``contiguous array'' question.}

% ------------------------------------------------------------------
\problem{Group anagrams}{Medium}

Group the words that are anagrams of one another.

\noindent\texttt{["eat","tea","tan","ate"]} gives
\texttt{[[eat,tea,ate],[tan]]}.

\idea{Give every word a canonical key that all its anagrams share. Sorting its
letters is the simplest such key.}

\begin{dia}
  \node[cell, minimum width=13mm] (w1) at (0,1.2) {eat};
  \node[cell, minimum width=13mm] (w2) at (0,0.5) {tea};
  \node[cell, minimum width=13mm] (w3) at (0,-0.2) {tan};
  \node[cell, minimum width=13mm] (w4) at (0,-0.9) {ate};
  \node[soft, minimum width=14mm] (k1) at (2.6,0.85) {aet};
  \node[soft, minimum width=14mm] (k2) at (2.6,-0.55) {ant};
  \node[note] at (2.6,1.5) {sorted key};
  \node[hot, minimum width=30mm] (g1) at (5.6,0.85) {eat, tea, ate};
  \node[hot, minimum width=30mm] (g2) at (5.6,-0.55) {tan};
  \draw[flow] (w1) -- (k1); \draw[flow] (w2) -- (k1); \draw[flow] (w4) -- (k1);
  \draw[flow] (w3) -- (k2);
  \draw[flow] (k1) -- (g1); \draw[flow] (k2) -- (g2);
\end{dia}

\textbf{How it works.} Two words are anagrams precisely when they have the
same multiset of letters. A multiset has no natural identity in Java, so you
manufacture one: sort the letters and the result is identical for all members
of a group and different for everything else. That string becomes the map
key, and \texttt{computeIfAbsent} does the grouping in one line.

A 26-slot count rendered as a string is a faster key --- it drops the
$\log k$ sorting factor --- but sorting is shorter to write and almost always
fast enough.

\solheading
\begin{lstlisting}[style=java]
public static List<List<String>> groupAnagrams(String[] words) {
    Map<String,List<String>> groups = new HashMap<>();
    for (String w : words) {
        char[] letters = w.toCharArray();
        Arrays.sort(letters);
        String key = new String(letters);
        groups.computeIfAbsent(key, k -> new ArrayList<>()).add(w);
    }
    return new ArrayList<>(groups.values());
}
\end{lstlisting}

\cost{$O(n\,k\log k)$ for $n$ words of length $k$.}
\trap{Using the word's \texttt{hashCode} or the sum of its characters as the
key. Both collide --- \texttt{"ad"} and \texttt{"bc"} have the same letter
sum --- and the bug only shows on inputs you did not try.}

% ------------------------------------------------------------------
\problem{Top k frequent elements}{Medium}

Return the $k$ most frequent values, most frequent first.

\idea{Count, then select. A min-heap capped at size $k$ keeps the selection at
$O(n\log k)$.}

\begin{dia}
  \node[note, anchor=east] at (-0.6,0) {counts};
  \node[soft, minimum width=12mm] at (0,0) {1:3};
  \node[soft, minimum width=12mm] at (1.3,0) {2:2};
  \node[soft, minimum width=12mm] at (2.6,0) {3:1};
  \node[soft, minimum width=12mm] at (3.9,0) {4:5};
  \node[note, anchor=east] at (-0.6,-1.4) {heap, $k=2$};
  \node[hot, minimum width=12mm] (h1) at (0,-1.4) {2:2};
  \node[hot, minimum width=12mm] (h2) at (1.3,-1.4) {4:5};
  \node[dead, minimum width=12mm] (h3) at (2.9,-1.4) {3:1};
  \node[note, anchor=west] at (3.6,-1.4) {evicted: smallest count};
  \draw[back] (h1) -- (h3);
  \caplabel{0.65}{-2.1}{the root of a min-heap is the \emph{least} frequent};
\end{dia}

\textbf{How it works.} Counting is the easy half. The selection is where the
choice lies: sorting all distinct values is $O(d\log d)$, while a heap capped
at $k$ is $O(d\log k)$ and, more importantly, needs only $k$ entries in
memory --- which is what makes it work on a stream.

The counter-intuitive part is that you keep the top $k$ with a \emph{min}-heap.
Push everything; whenever the heap exceeds $k$, discard its root, which is the
least frequent thing currently held. What survives is the $k$ most frequent.
Because the root is the smallest, draining the heap yields ascending
frequency, so the output array is filled backwards.

\solheading
\begin{lstlisting}[style=java]
public static int[] topKFrequent(int[] a, int k) {
    Map<Integer,Integer> freq = new HashMap<>();
    for (int x : a) freq.merge(x, 1, Integer::sum);

    PriorityQueue<Map.Entry<Integer,Integer>> heap =
        new PriorityQueue<>(Map.Entry.comparingByValue());   // min-heap
    for (Map.Entry<Integer,Integer> e : freq.entrySet()) {
        heap.offer(e);
        if (heap.size() > k) heap.poll();     // drop the least frequent
    }
    int[] out = new int[heap.size()];
    for (int i = out.length - 1; i >= 0; i--) out[i] = heap.poll().getKey();
    return out;
}
\end{lstlisting}

\cost{$O(n\log k)$ time, $O(n)$ space for the counts.}
\trap{Filling the output forwards. The heap drains smallest-first, so a
forward fill returns the $k$ most frequent in exactly the wrong order.}

% ------------------------------------------------------------------
\problem{Subarray sum equals k}{Medium}

Count the contiguous subarrays summing to exactly $k$. Values may be negative.

\idea{Prefix sums. A window sums to $k$ whenever an earlier prefix equals the
current prefix minus $k$.}

\begin{dia}
  \cellrowst{a}{0.9}{cell/1,cell/2,cell/3,cell/-2}
  \node[note, anchor=east] at (-0.5,0.9) {input};
  \cellrowst{p}{0}{soft/0,soft/1,soft/3,hot/6,soft/4}
  \node[note, anchor=east] at (-0.5,0) {prefix};
  \node[idx] at (-0.0,-0.55) {$\varnothing$};
  \draw[span] (0.4,-1.0) -- (2.0,-1.0);
  \caplabel{1.2}{-1.5}{$6-1=5$: the window summing to 5 is $2+3$};
  \draw[back] (2.4,0.45) to[out=150,in=30] (0.8,0.45);
\end{dia}

\textbf{How it works.} Let $P_i$ be the sum of the first $i$ elements. The sum
of the window $(j,i]$ is $P_i - P_j$. Asking for a window summing to $k$ is
therefore asking: how many earlier prefixes equal $P_i - k$? Keep a map from
prefix value to how many times it has occurred, and the answer accumulates in
one pass.

The map is seeded with $0 \mapsto 1$ for the empty prefix, which is what lets
a window starting at index 0 be counted.

It is worth knowing why a sliding window fails here. Windows work when
growing the window can only increase the sum, so you can shrink from the left
when you overshoot. Negative values destroy that monotonicity; prefix sums do
not care.

\solheading
\begin{lstlisting}[style=java]
public static int subarraysSummingTo(int[] a, int k) {
    Map<Long,Integer> counts = new HashMap<>();
    counts.put(0L, 1);                     // the empty prefix
    long prefix = 0;
    int total = 0;
    for (int x : a) {
        prefix += x;
        total += counts.getOrDefault(prefix - k, 0);
        counts.merge(prefix, 1, Integer::sum);
    }
    return total;
}
\end{lstlisting}

\cost{$O(n)$ time, $O(n)$ space.}
\trap{Reaching for a sliding window because the problem says ``contiguous''.
Check for negative values first --- they are the tell.}

% ------------------------------------------------------------------
\problem{First duplicate}{Easy}

Return the first value that appears twice, scanning left to right, or $-1$.

\idea{A \texttt{HashSet}. \texttt{add} returns \texttt{false} when the element
was already there, so the check and the insert are a single call.}

\begin{dia}
  \cellrowst{a}{0}{dead/2,dead/1,dead/3,hot/1,cell/5}
  \idxrow{-0.62}{4}
  \node[soft, minimum width=9mm] at (0,-1.5) {2};
  \node[soft, minimum width=9mm] at (1.0,-1.5) {1};
  \node[soft, minimum width=9mm] at (2.0,-1.5) {3};
  \node[note, anchor=east] at (-0.5,-1.5) {set};
  \draw[flow] (2.4,-0.95) -- (1.1,-1.15);
  \caplabel{3.6}{-1.5}{\texttt{add(1)} returns false $\Rightarrow$ answer};
\end{dia}

\textbf{How it works.} \texttt{Set.add} is specified to return whether the set
changed. That single boolean replaces a \texttt{contains} followed by an
\texttt{add}, halving the number of hash lookups and removing the chance of
writing one without the other.

Be sure which question is being asked. This returns the value whose
\emph{second} occurrence comes earliest. ``The duplicate whose first
occurrence is earliest'' is a different answer and needs a second pass.

\solheading
\begin{lstlisting}[style=java]
public static int firstDuplicate(int[] a) {
    Set<Integer> seen = new HashSet<>();
    for (int x : a) if (!seen.add(x)) return x;
    return -1;
}
\end{lstlisting}

\cost{$O(n)$ time, $O(n)$ space.}
\trap{Misreading which ``first'' is wanted. Also, $-1$ as a sentinel is wrong
if $-1$ can itself be a duplicate value --- return \texttt{Integer} and
\texttt{null} if the statement allows negatives.}

% ------------------------------------------------------------------
\problem{Intersection of two arrays}{Easy}

Return the values present in both arrays, each once.

\idea{Put one array in a set, filter the other against it, and collect the hits
into a second set so duplicates collapse.}

\begin{dia}
  \draw[edge, fill=dsoft, fill opacity=0.6] (0,0) circle (1.1);
  \draw[edge, fill=dwarm, fill opacity=0.6] (1.5,0) circle (1.1);
  \node[note] at (-0.55,0) {2, 9};
  \node[note] at (0.75,0) {4, 5};
  \node[note] at (2.05,0) {7};
  \node[note] at (-0.55,1.3) {a};
  \node[note] at (2.05,1.3) {b};
  \node[note, anchor=west] at (3.2,0) {result: \texttt{[4, 5]}};
\end{dia}

\textbf{How it works.} The first set turns ``is this value in \texttt{a}'' from
a scan into an $O(1)$ question, which is the whole saving over the nested
loop. The second set exists for a different reason: if \texttt{b} contains a
shared value twice, a plain list would emit it twice.

\texttt{LinkedHashSet} is used rather than \texttt{HashSet} so the output
order is the order of first appearance in \texttt{b}. Plain \texttt{HashSet}
leaves the order unspecified, and graders usually compare arrays exactly.

\solheading
\begin{lstlisting}[style=java]
public static int[] intersection(int[] a, int[] b) {
    Set<Integer> inA = new HashSet<>();
    for (int x : a) inA.add(x);
    Set<Integer> both = new LinkedHashSet<>();
    for (int x : b) if (inA.contains(x)) both.add(x);
    int[] out = new int[both.size()];
    int i = 0;
    for (int x : both) out[i++] = x;
    return out;
}
\end{lstlisting}

\cost{$O(n+m)$ time, $O(n)$ space.}
\trap{Collecting into a \texttt{List}, which repeats values that recur in
\texttt{b}.}

% ------------------------------------------------------------------
\problem{Isomorphic strings}{Medium}

Two strings are isomorphic if the characters of one map onto the characters of
the other consistently and one-to-one. \texttt{"egg"} and \texttt{"add"} are;
\texttt{"foo"} and \texttt{"bar"} are not.

\idea{Two maps, one in each direction. One map alone permits two different
characters to land on the same target.}

\begin{dia}
  \node[cell] (e1) at (0,0.5) {f};
  \node[cell] (e2) at (0.8,0.5) {o};
  \node[cell] (e3) at (1.6,0.5) {o};
  \node[cell] (a1) at (0,-1.1) {b};
  \node[cell] (a2) at (0.8,-1.1) {a};
  \node[hot]  (a3) at (1.6,-1.1) {r};
  \draw[flow] (e1) -- (a1);
  \draw[flow] (e2) -- (a2);
  \draw[flow] (e3) -- (a3);
  \node[note, anchor=west] at (2.4,-0.3) {\texttt{o} already maps to \texttt{a},};
  \node[note, anchor=west] at (2.4,-0.8) {so \texttt{o}$\to$\texttt{r} contradicts it};
\end{dia}

\textbf{How it works.} Consistency in one direction says: the same source
character always produces the same target. One-to-one adds the reverse: no two
source characters share a target. \texttt{"ab"} and \texttt{"cc"} satisfies
the first and fails the second, which is why a single map is not enough.

\texttt{putIfAbsent} is doing double duty: it records the mapping if it is new
and returns the existing one if it is not, so the check and the write are one
call per direction.

\solheading
\begin{lstlisting}[style=java]
public static boolean isIsomorphic(String s, String t) {
    if (s == null || t == null || s.length() != t.length()) return false;
    Map<Character,Character> forward = new HashMap<>();
    Map<Character,Character> backward = new HashMap<>();
    for (int i = 0; i < s.length(); i++) {
        char a = s.charAt(i), b = t.charAt(i);
        Character f = forward.putIfAbsent(a, b);
        Character g = backward.putIfAbsent(b, a);
        if ((f != null && f != b) || (g != null && g != a)) return false;
    }
    return true;
}
\end{lstlisting}

\cost{$O(n)$ time, $O(1)$ space for a bounded alphabet.}
\trap{Comparing the boxed \texttt{Character} values with \texttt{!=} works here
only because the cache covers ASCII. With \texttt{Integer} keys above 127 the
identical code fails silently --- make unboxing a habit.}
